% =====================================================================
%  CHƯƠNG 7. HỆ THỐNG QUẢN LÝ BÃI GIỮ XE       | Phụ trách: Đức
%  Ngân sách: mục tiêu 6 trang. Viết high-level, không đi sâu vào mã nguồn.
%  Nội dung dựa trên thiết kế backend (FastAPI, PostgreSQL) và dashboard (React).
% =====================================================================
\chapter{Hệ Thống Quản Lý Bãi Giữ Xe}
\label{ch:hethong}

% -------------------------------------------------------------------
\section{Kiến Trúc Tổng Thể}
\label{sec:kientruc}
% -------------------------------------------------------------------

Hệ thống tổ chức theo ba lớp: thu nhận ảnh tại cổng, backend xử lý nghiệp vụ, và cơ sở dữ
liệu. Điểm khác biệt là hai đường thu nhận ảnh song song cùng đổ về một backend duy nhất và
dùng chung một bộ suy luận (Hình~\ref{fig:kientruc}). Đường thứ nhất là kiosk trên máy tính:
giao diện cổng chụp ảnh từ webcam, gửi lên backend để suy luận ngay trong tiến trình và trả
kết quả. Đường thứ hai là cổng tự động trên thiết bị biên: worker chạy suy luận tại chỗ rồi
gửi kết quả kèm ảnh về backend, backend không suy luận lại mà chỉ ghi nhận và xử lý nghiệp
vụ. Cả hai đường dùng chung một pipeline nhận diện, nên kết quả nhất quán bất kể ảnh được xử
lý ở đâu.

\begin{figure}[htbp]
\centering
\includegraphics[width=0.9\textwidth]{uml-kientruc.png}
\caption{Kiến trúc tổng thể hệ thống: hai đường thu nhận ảnh cùng đổ về một backend.}
\label{fig:kientruc}

\end{figure}

Cùng một bộ mã chạy trên hai cấu hình triển khai, chọn qua cấu hình mà không sửa mã
(Bảng~\ref{tab:cauhinh-trienkhai}): trên máy tính, các thành phần chạy trong container với
PostgreSQL; trên Raspberry Pi~5, backend và giao diện chạy trực tiếp trên hệ điều hành với
SQLite và tăng số luồng suy luận để tận dụng bốn nhân.

\begin{table}[htbp]
\centering
\caption{Hai cấu hình triển khai của hệ thống.}
\label{tab:cauhinh-trienkhai}

\vspace{4pt}
\begin{tabular}{lll}
\toprule
\textbf{Thành phần} & \textbf{Máy tính} & \textbf{Raspberry Pi 5} \\
\midrule
Backend & Container FastAPI & FastAPI chạy trực tiếp \\
Cơ sở dữ liệu & PostgreSQL & SQLite \\
Phát hiện biển số & YOLOv8n & YOLOv8n \\
Đọc ký tự, loại xe, kiểu dáng & ONNX Runtime & ONNX Runtime \\
Số luồng mỗi mô hình & 1 & 4 \\
\bottomrule
\end{tabular}
\end{table}

Pipeline nạp các mô hình một lần khi khởi động rồi chạy trọn chuỗi nhận diện
(Mục~\ref{sec:donggop}) trên mỗi khung ảnh; nhờ tách khỏi lớp vận chuyển, cùng một pipeline
phục vụ cả kiosk lẫn worker biên.

% -------------------------------------------------------------------
\section{Cơ Sở Dữ Liệu}
\label{sec:csdl}
% -------------------------------------------------------------------

Cơ sở dữ liệu tổ chức quanh tám thực thể trung tâm (Bảng~\ref{tab:bang-chinh}), với sơ đồ
quan hệ ở Hình~\ref{fig:uml-er}. Ngoài ra còn các bảng hỗ trợ cho nhóm xe và bảng giá, vé
tháng, danh sách trắng và đen, sự cố, thiết bị cổng và cấu hình bãi.

\begin{table}[htbp]
\centering
\caption{Các thực thể trung tâm trong cơ sở dữ liệu.}
\label{tab:bang-chinh}

\vspace{4pt}
\begin{tabular}{lp{9.5cm}}
\toprule
\textbf{Thực thể} & \textbf{Vai trò} \\
\midrule
Phiên gửi xe & Vòng đời một lượt gửi xe: trạng thái, phí chốt và bản ghi quy tắc giá tại thời điểm chốt \\
Lượt đọc biển & Một lần đọc biển tại cổng (vào hoặc ra), lưu kết quả nhận diện và trạng thái duyệt \\
Thanh toán & Giao dịch thu tiền, hoàn trả hoặc điều chỉnh theo ca trực \\
Quy tắc giá & Bảng giá theo nhóm xe: chế độ tính, đơn giá, khối thời gian, thời gian miễn phí, mức trần ngày \\
Ảnh bằng chứng & Ảnh toàn cảnh và ảnh cắt biển; trạng thái mã hóa và thời hạn xóa \\
Ca trực & Ca làm việc của nhân viên, gắn với các giao dịch thu tiền \\
Người dùng & Tài khoản và vai (nhân viên, quản lý, quản trị) \\
Nhật ký kiểm toán & Ghi lại các thao tác nhạy cảm để đối soát sau sự cố \\
\bottomrule
\end{tabular}
\end{table}

\begin{figure}[htbp]
\centering
\includegraphics[width=0.9\textwidth]{uml-er.png}
\caption{Sơ đồ quan hệ thực thể của các bảng trung tâm.}
\label{fig:uml-er}

\end{figure}

\textbf{Bảo mật dữ liệu cá nhân.} Biển số và ảnh phương tiện là dữ liệu cá nhân theo Luật Bảo
vệ dữ liệu cá nhân~\autocite{vn2025lutbvdlcn}. Mỗi chuỗi biển được lưu đồng thời dạng băm (để
đối chiếu) và dạng mã hóa (để hiển thị khi có quyền); ảnh bằng chứng được mã hóa và chỉ truy
cập qua endpoint có kiểm tra quyền; một tiến trình nền xóa dữ liệu quá hạn tính từ lúc xe ra.
Đối chiếu đầy đủ với nghĩa vụ theo luật ở Phụ lục~\ref{app:baomat}.

% -------------------------------------------------------------------
\section{Logic Vào Ra và Chống Gian Lận}
\label{sec:logic-vaora}
% -------------------------------------------------------------------

Mỗi phiên gửi xe đi qua các trạng thái: đang trong bãi, chờ xử lý thủ công (khi không đọc
được biển), hoàn thành, và tranh chấp. Khi xe vào, hệ thống tạo phiên mới và cảnh báo nếu
biển trùng một phiên đang trong bãi hoặc nằm trong danh sách đen.

Khi xe ra, hệ thống đối chiếu biển theo ba cấp (Hình~\ref{fig:uml-seq}). Trước hết là khớp
chính xác qua giá trị băm; nếu không khớp, hệ thống so khoảng cách chỉnh sửa giữa chuỗi biển
ra và các phiên đang trong bãi cùng nhóm xe, chấp nhận khi lệch không quá một ký tự (xử lý
sai sót nhỏ của OCR); nếu vẫn không rõ ràng, hệ thống đưa danh sách ứng viên cho nhân viên
chọn hoặc chuyển phiên sang trạng thái tranh chấp. Mỗi lượt đọc đều gắn ảnh toàn cảnh và ảnh
cắt biển làm bằng chứng khi cần đối soát; trường hợp mất vé được xử lý thành phiên riêng có
phí phạt và ghi nhật ký.

\begin{figure}[htbp]
\centering
\includegraphics[width=0.9\textwidth]{uml-seq.png}
\caption{Sơ đồ tuần tự một lượt xe từ vào đến ra: thu nhận ảnh, tạo phiên, đối chiếu biển, tính phí và ghi thanh toán.}
\label{fig:uml-seq}

\end{figure}

% -------------------------------------------------------------------
\section{Tính Phí Tự Động}
\label{sec:tinphi}
% -------------------------------------------------------------------

Mỗi nhóm xe (xe máy, ô tô con, xe tải, xe khách) có một quy tắc giá đang kích hoạt. Hệ thống
hỗ trợ hai chế độ: chế độ phẳng thu cố định mỗi lượt, và chế độ theo khối thời gian nhân đơn
giá với số khối làm tròn lên, có thể chặn tại mức trần theo ngày. Khoảng thời gian miễn phí
đầu (grace period) cho phép ra trong ít phút mà không mất phí. Khi phiên kết thúc, toàn bộ
quy tắc được chốt kèm phiên tại thời điểm đó, nên thay đổi bảng giá về sau không ảnh hưởng
phí đã tính của các phiên cũ và vẫn đối soát được. Xe thuộc danh sách miễn phí được ghi rõ
lý do miễn để phân biệt với phí thật sự bằng không.

% -------------------------------------------------------------------
\section{Dashboard Quản Lý và Thống Kê}
\label{sec:dashboard}
% -------------------------------------------------------------------

Hệ thống phân ba vai người dùng cùng một tác nhân thiết bị biên
(Bảng~\ref{tab:vaitro}, sơ đồ use case ở Hình~\ref{fig:uml-usecase}). Nhân viên xử lý nghiệp
vụ cổng; quản lý thêm quyền cấu hình bảng giá, tra cứu thống kê và xuất dữ liệu; quản trị
quản lý tài khoản. Thiết bị biên là tác nhân riêng, xác thực bằng khóa riêng thay vì tài
khoản người dùng.

\begin{table}[htbp]
\centering
\caption{Các tác nhân và quyền chính trong hệ thống.}
\label{tab:vaitro}

\vspace{4pt}
\begin{tabular}{lp{9.5cm}}
\toprule
\textbf{Tác nhân} & \textbf{Quyền chính} \\
\midrule
Nhân viên & Chụp và nhận diện, xác nhận vào ra, thu tiền, sửa biển, mở và chốt ca, tra cứu phiên \\
Quản lý & Toàn bộ quyền nhân viên; bảng giá, nhóm xe, vé tháng, danh sách trắng và đen, giải quyết tranh chấp, thống kê và xuất CSV, tạo tài khoản nhân viên \\
Quản trị & Toàn bộ quyền quản lý; quản lý tài khoản vai quản lý và quản trị \\
Thiết bị biên & Gửi kết quả nhận diện kèm ảnh và tín hiệu hoạt động; xác thực bằng khóa riêng \\
\bottomrule
\end{tabular}
\end{table}

\begin{figure}[htbp]
\centering
\includegraphics[width=0.9\textwidth]{uml-usecase.png}
\caption{Sơ đồ use case: nhân viên thực hiện nghiệp vụ cổng, quản lý thêm quyền cấu hình và thống kê.}
\label{fig:uml-usecase}

\end{figure}

Giao diện gồm bốn nhóm màn hình chính. Màn cổng chia khung vào và khung ra, hiển thị kết quả
nhận diện theo thời gian thực: biển số, loại xe, màu nền, ảnh biển đã xử lý và phí đề xuất
(Hình~\ref{fig:dash-gate-in}); khi xe ra, khung ra đặt ảnh lúc vào cạnh ảnh lúc ra để đối
chiếu trước khi thu phí (Hình~\ref{fig:dash-gate-out}). Trang tra cứu phiên lọc theo nhóm xe,
thời gian, trạng thái và cách khớp biển (Hình~\ref{fig:dash-sessions}); trang chi tiết giữ
ảnh bằng chứng của cả hai lượt cùng phí và quy tắc giá lúc chốt
(Hình~\ref{fig:dash-detail}). Trang thống kê tổng hợp số xe vào ra và doanh thu theo ngày,
kèm xuất CSV (Hình~\ref{fig:dash-stats}). Trang cấu hình quản lý bảng giá, tài khoản và công
tắc các bước AI, đồng thời nêu rõ chính sách tự xóa dữ liệu phiên sau thời hạn quy định
(Hình~\ref{fig:dash-config}).

\begin{figure}[htbp]
\centering
\includegraphics[width=0.92\textwidth]{dashboard-gate-in.png}
\caption{Màn cổng, khung vào: kết quả nhận diện thời gian thực gồm biển số, loại xe, màu nền, ảnh biển đã xử lý và phí đề xuất.}
\label{fig:dash-gate-in}

\end{figure}

\begin{figure}[htbp]
\centering
\includegraphics[width=0.92\textwidth]{dashboard-gate-out.png}
\caption{Khung ra đối chiếu ảnh lúc vào và lúc ra của cùng một xe trước khi thu phí.}
\label{fig:dash-gate-out}

\end{figure}

\begin{figure}[htbp]
\centering
\includegraphics[width=0.95\textwidth]{dashboard-sessions.png}
\caption{Trang tra cứu phiên: bộ lọc và danh sách phiên với biển số, thời gian vào ra, phí, cách khớp và trạng thái.}
\label{fig:dash-sessions}

\end{figure}

\begin{figure}[htbp]
\centering
\includegraphics[width=0.92\textwidth]{dashboard-session-detail.png}
\caption{Chi tiết một phiên gửi xe: thông tin phiên, phân loại và ảnh bằng chứng lúc vào, lúc ra.}
\label{fig:dash-detail}

\end{figure}

\begin{figure}[htbp]
\centering
\includegraphics[width=0.95\textwidth]{dashboard-stats-new.png}
\caption{Trang thống kê: thẻ tổng hợp, biểu đồ lưu lượng và doanh thu theo ngày, nút xuất CSV.}
\label{fig:dash-stats}

\end{figure}

\begin{figure}[htbp]
\centering
\includegraphics[width=0.95\textwidth]{dashboard-config.png}
\caption{Trang cấu hình bảng giá theo nhóm xe, kèm ghi chú tự xóa dữ liệu phiên theo Luật Bảo vệ dữ liệu cá nhân.}
\label{fig:dash-config}

\end{figure}
